\section{Representaciones de Galois y compatibilidad local}
\label{sec:galois}

Las congruencias automorfas de las notas anteriores deben convertirse en
información aritmética sobre representaciones de Galois. Para hacerlo
no basta producir matrices con los polinomios de Frobenius esperados:
hay que obtener la dimensión correcta, transportar las condiciones
locales y conservar las normalizaciones que después intervienen en la
cohomología. Esta nota desarrolla los mecanismos de extracción y
descenso que realizan esas tareas, una vez recibidas las entradas
geométricas y automorfas.

\fuente{A, §9, pp.~42--46; A93, §10, pp.~77--85.
Se conserva A como nombre del ejemplar de 55 páginas.
Las entradas \texttt{Row0b} y \texttt{RecTExists} se explicitan;
las pruebas de álgebra de representaciones y sus consecuencias se
desarrollan aquí. El cotejo del código se delimita en
§\ref{subsec:gal-lean}. La mejora de la conclusión del pegado en
la Proposición~\ref{prop:gal-todos} es una deducción adicional de
estas notas, no el enunciado literal de A93 ni de Lean.}

\subsection{La función de esta entrada en el cuaderno}

Escribimos $G_E=\operatorname{Gal}(\overline E/E)$ para el grupo de
Galois absoluto de un cuerpo de números $E$. Una representación
$\ell$-ádica de dimensión $a$ es un homomorfismo continuo
\[
 \rho:G_E\longrightarrow
       \operatorname{GL}(V),\qquad
 \dim_{\overline\Q_\ell}V=a.
\]
Un subespacio es invariante si lo preservan todas las matrices
$\rho(g)$. La representación es \emph{irreducible} si no tiene
subespacios invariantes propios no nulos, y \emph{semisimple} si
es suma directa de irreducibles. Una isomorfía de representaciones
es una transformación lineal invertible que entrelaza todas
las acciones.

Para un lugar finito $v$ de $E$, una elección de lugar sobre $v$
en $\overline E$ determina un grupo de descomposición $D_v$ y su
subgrupo de inercia $I_v$. La representación es no ramificada
en $v$ cuando $\rho(I_v)=1$. En ese caso tiene sentido su polinomio
característico de Frobenius. La elección de un lugar conjugado
cambia las matrices por conjugación, de modo que no cambia ese
polinomio.

La representación automorfa $\pi$ pertenece a
$\operatorname{GL}_a(\mathbb A_E)$, donde $\mathbb A_E$ es el anillo
de adeles. Su espacio no tiene por qué ser finito dimensional.
El entero $a$ es el rango del grupo lineal y será la dimensión de
la representación de Galois asociada. La regularidad algebraica
prescribe datos infinitesimales en los lugares complejos; esos datos
se convierten, con la normalización adecuada, en pesos de Hodge--Tate.

El recorrido de la construcción es el siguiente:
\begin{enumerate}
 \item La geometría proporciona representaciones auxiliares,
       posiblemente de dimensión mayor que $a$.
 \item Los levantamientos theta identifican los sumandos adicionales
       como listas de caracteres.
 \item La extracción produce una representación genuina de dimensión $a$
       y conserva las propiedades locales de los sumandos.
 \item El pegado cuadrático retira restricciones sobre el cuerpo base.
 \item Las comparaciones en lugares que se descomponen recuperan
       las condiciones locales sobre el cuerpo original.
\end{enumerate}
La existencia de representaciones automorfas de Galois tiene una
literatura previa, como \cite{chenevier-harris}. Aquí interesa
la forma precisa de la entrada que A93 usa para sus estimaciones,
sin identificarla con todos los resultados posibles de esa teoría.

\subsection{La entrada \texttt{Row0b}: hipótesis y pesos}
\label{subsec:gal-row}

Sea $E$ un cuerpo CM, $F$ su subcuerpo totalmente real maximal
y $c$ la conjugación compleja de $E/F$. Fijamos $a\ge1$,
una representación automorfa cuspidal regular algebraica $\pi$
de $\operatorname{GL}_a(\mathbb A_E)$ y una identificación de cuerpos
$\iota:\overline\Q_\ell\simeq\C$. La polarización utilizada por la
fuente incluye un carácter de Hecke $\chi$ de $F$ y las dos condiciones
\begin{equation}
 \pi^c\simeq\pi^\vee\otimes(\chi\circ N_{E/F}),
 \qquad \chi_v(-1)=(-1)^a
       \quad\text{para cada lugar real }v\text{ de }F.
 \label{eq:gal-polarizacion}
\end{equation}
La condición de signo forma parte de la hipótesis;
no se omite al abreviar la dualidad conjugada.
En el código ambas aparecen en \texttt{PolarizedPair}.

La normalización algebraica de la componente local se denota
$\pi_v^{\mathrm{alg}}$. La entrada afirma la existencia de una
representación continua semisimple $\rho_{\pi,\ell}$ de dimensión $a$,
un entero $w$ y multiconjuntos integrales $H_\tau$ tales que
\begin{equation}
 \begin{split}
 H_\tau&=\{(a-1)/2-z:
                  z\in\operatorname{Inf}_{\iota\tau}(\pi)\},\\
 H_{\tau c}&=\{w-h:h\in H_\tau\}.
 \end{split}
 \label{eq:gal-pesos}
\end{equation}
Aquí $\tau:E\hookrightarrow\overline\Q_\ell$,
$\operatorname{Inf}_{\iota\tau}(\pi)$ es el multiconjunto de
parámetros del carácter infinitesimal en la convención de la fuente
y $\tau c=\tau\circ c$. Todas las multiplicidades se conservan.

Un multiconjunto registra cuántas veces aparece cada elemento.
Por ejemplo, si $u\ne v$, los multiconjuntos
$\{u,u,v\}$ y $\{u,v,v\}$ tienen la misma cardinalidad y el mismo
conjunto de valores, pero no son iguales. Los operadores diagonales
correspondientes son aniquilados por $(T-u)(T-v)$, aunque sus
polinomios característicos son, respectivamente,
\begin{equation}
 (T-u)^2(T-v),\qquad (T-u)(T-v)^2.
 \label{eq:gal-multiplicidades}
\end{equation}
Una relación de anulación no determina por sí sola los pesos
con multiplicidad, ni la multiplicidad de los eigenvalores de Frobenius.

En esta nota los pesos de Hodge son los índices $h$ de los saltos
$\operatorname{gr}^h D_{\mathrm{dR}}$; con esta convención el carácter
ciclotómico tiene peso $-1$. Los polinomios no ramificados se
comparan con la misma elección de Frobenius y el mismo giro
algebraico en ambos lados. Cambiar una sola de esas convenciones
alteraría \eqref{eq:gal-pesos} o los factores de las comparaciones.

El signo del peso ciclotómico puede comprobarse con el periodo
$t$ de filtración $1$, que satisface $g(t)=\chi_\ell(g)t$:
si $e$ genera la representación ciclotómica, $t^{-1}e$ es
invariante y tiene salto de filtración $-1$.

\subsection{Dos niveles de compatibilidad local}

Un parámetro de Weil--Deligne es un par $y=(r_y,N_y)$:
$r_y$ es una representación del grupo de Weil con imagen finita
en inercia y $N_y$ es un operador nilpotente compatible con
la norma del grupo de Weil. La monodromía es el operador $N_y$.
Sea
\[
 y_v=\operatorname{rec}^{\mathrm{aux}}_v(\pi_v^{\mathrm{alg}})
\]
la asignación local auxiliar de A93. Si $v\nmid\ell$, denotamos
por $x_v$ el parámetro asociado a $\rho_{\pi,\ell}$.

Llamamos \emph{abeliano} a un parámetro cuando su representación
de Weil factoriza por la abelianización. La condición mansa
de multiplicidad exige que la inercia salvaje actúe trivialmente y que,
para un generador $\gamma$ de la inercia mansa,
\begin{equation}
 \operatorname{mult}_{r_y(\gamma)}(z)\ge2
 \quad\Longrightarrow\quad z^{q_v-1}=1,
 \label{eq:gal-mansa}
\end{equation}
donde $q_v$ es la cardinalidad del cuerpo residual.
La segunda hipótesis de la rama mansa es
\begin{equation}
 E_v\simeq\Q_p,\qquad a=1\ \text{o}\ p>16a^2.
 \label{eq:gal-prime-field}
\end{equation}
Una extensión propia de $\Q_p$ no satisface esta hipótesis
por el solo hecho de tener $p$ grande.

La parte de \texttt{Row0b} que utiliza el cuaderno puede organizarse así:
\begin{center}
\small
\begin{tabularx}{\textwidth}{@{}p{0.20\textwidth}YY@{}}
\toprule
Lugar & Hipótesis local adicional & Conclusión recibida\\
\midrule
$v\nmid\ell$ &
$N_{y_v}=0$ y $y_v$ abeliano &
Igualdad de parámetros $x_v=y_v$.\\
$v\nmid\ell$ &
$N_{y_v}=0$, \eqref{eq:gal-mansa} y \eqref{eq:gal-prime-field} &
$N_{x_v}=0$ e igualdad de polinomios característicos en inercia.\\
$v\mid\ell$ &
Componente esférica, en una realización compatible &
Cristalinidad, pesos $H_\tau$ y comparación no ramificada de Frobenius.\\
$v\mid\ell$ &
$N_{y_v}=0$ y una de las dos ramas anteriores, en una realización compatible &
Propiedad de de Rham, semiestabilidad potencial, pesos $H_\tau$
y la comparación de Weil--Deligne de la rama correspondiente.\\
\bottomrule
\end{tabularx}
\end{center}
La igualdad de dimensiones también forma parte de la entrada.
Una misma $\rho_{\pi,\ell}$ sirve para toda asignación local
admitida por las leyes de la fuente y para toda realización
compatible. En un factor local, los pesos $H_\tau$ corresponden
a las inmersiones globales que inducen ese lugar mediante el puente.
La tabla conserva la diferencia entre igualdad del parámetro completo
e igualdad de polinomios en inercia. Su formulación procede de
A93, Proposición~10.1, pp.~77--78, y de
\texttt{WeightsConjBis.lean}; aquí se recibe como entrada.

La monodromía no se deduce de los polinomios de inercia.
Para verlo, elijamos un elemento de Frobenius $\Phi$ y pongamos
$c_\Phi=|\Phi|$. Considérese el parámetro no ramificado de Weil
con $r(\Phi)=\operatorname{diag}(c_\Phi,1)$.
Tanto $N=0$ como
\begin{equation}
 N=\begin{pmatrix}0&1\\0&0\end{pmatrix}
 \quad\text{satisfacen}\quad
 r(\Phi)Nr(\Phi)^{-1}=c_\Phi N.
 \label{eq:gal-monodromia}
\end{equation}
Ambos tienen polinomio $(T-1)^2$ en cada elemento de inercia,
y el segundo tiene monodromía no nula.
Por ello la hipótesis $N_{y_v}=0$ y la conclusión $N_{x_v}=0$
tienen que acompañar la comparación mansa.

\subsection{Realizaciones locales con testigos}

La expresión «en toda realización compatible» tiene contenido
matemático. Para $v\mid\ell$ se elige un anillo de valoración
$A\subset\overline E$ sobre $v$ y una inmersión
$j:\overline E\hookrightarrow\overline\Q_\ell$ que induzca
esa valoración:
\begin{equation}
 x\in A\quad\Longleftrightarrow\quad |j(x)|_\ell\le1.
 \label{eq:gal-induces}
\end{equation}
El cuerpo local correspondiente es $\Q_\ell(j(E))$.
La identificación del grupo local debe respetar la acción:
\begin{equation}
 e(g)(j(x))=j(g(x)).
 \label{eq:gal-bridge}
\end{equation}
No basta escoger por separado una inmersión y unos anillos
de periodos y suponer que corresponden al mismo lugar.

A93, p.~78, y las declaraciones
\texttt{exists\_locPrime\_tied} y \texttt{exists\_locBridge}
aportan testigos de existencia. La condición \texttt{LocTied}
liga los datos locales a los anillos de periodos cristalinos
y de de Rham; \texttt{Induces} expresa
\eqref{eq:gal-induces}. El puente se construye tras ampliar
suficientemente el cuerpo finito de coeficientes.
Así la cuantificación universal sobre realizaciones compatibles
puede instanciarse en cada lugar sobre $\ell$.

La lógica que se necesita es
\begin{equation}
 \exists\rho\quad\forall v\mid\ell\quad
 \left[
   \begin{array}{l}
    \text{existen datos locales compatibles en }v,\\
    \text{la conclusión vale para todo dato compatible.}
   \end{array}
 \right]
 \label{eq:gal-cuantificadores-locales}
\end{equation}
La existencia de datos y la validez universal de la comparación
son afirmaciones distintas. Esta distinción será importante al
transportar propiedades a un sumando de $\rho$.

\subsection{El alcance de la asignación local auxiliar}
\label{subsec:gal-rect}

La fuente llama \emph{buena} a una supercuspidal si es un carácter
de $\operatorname{GL}_1$, o si el cuerpo local es $\Q_p$ y
es un giro por un carácter de una supercuspidal de profundidad cero.
Una representación irreducible admisible es buena cuando todos
los miembros de su soporte supercuspidal lo son.
«Profundidad cero» se expresa mediante vectores fijos por el primer
subgrupo principal de congruencia del compacto maximal elegido.

La asignación cuya existencia expresa \texttt{RecTExists}
conserva la dimensión, coincide
con reciprocidad local en rango uno, respeta la descripción por
segmentos y los giros reales no ramificados, y envía las
supercuspidales a parámetros irreducibles. También prescribe
el valor absoluto del determinante. En la clase buena se añaden
las leyes de giro por caracteres arbitrarios, dualidad,
determinante exacto e inyectividad frente a un parámetro bueno,
junto con las comparaciones de tipos de profundidad cero.
Estas últimas leyes no se afirman sin esa restricción.

La distinción permite completar una asignación sobre todas las
clases locales sin presentar esa completación como una nueva
demostración de toda la correspondencia local de Langlands.
Las leyes que utiliza la aplicación deben verificarse en la
clase a la que pertenece su componente local.
El enunciado de A93 es la Proposición~10.6, p.~85;
el código separa \texttt{IsRecTOfIsRec} de
\texttt{LocalCentralUnits/RecTExists}.

\begin{lema}[Órbitas de inercia y multiplicidad]
\label{lem:gal-orbitas}
Supongamos que la inercia salvaje actúa trivialmente.
La condición \eqref{eq:gal-mansa} equivale a que toda órbita
de eigenvalores bajo $z\mapsto z^{q_v}$ de longitud mayor que uno
aparezca con multiplicidad uno.
Las órbitas fijas pueden tener multiplicidades mayores.
\end{lema}

\begin{proof}
La imagen finita de inercia es semisimple en característica cero.
La relación entre Frobenius aritmético e inercia permuta los
eigenvalores por $z\mapsto z^{q_v}$ y conserva sus multiplicidades.
El uso de Frobenius geométrico invierte esta permutación y deja
las mismas órbitas. Un eigenvalor es fijo precisamente cuando
$z^{q_v-1}=1$. Si una órbita no fija tuviese multiplicidad al menos
dos, violaría \eqref{eq:gal-mansa}. Recíprocamente, si las órbitas
no fijas son simples, todo eigenvalor repetido es fijo.
\end{proof}

Por ejemplo, una órbita regular de caracteres de un cuerpo finito
da eigenvalores distintos en el parámetro inducido de profundidad
cero. La suma de varios caracteres de inercia triviales también
satisface \eqref{eq:gal-mansa}, aunque tiene raíces repetidas.
La condición no equivale a que el polinomio sea libre de cuadrados.

\subsection{La entrada geométrica y el levantamiento theta}

La construcción geométrica de A93, pp.~78--79, parte de datos
unitarios compactos y sistemas propios de Hecke en secciones
coherentes. Los invariantes de Hasse permiten mover pesos para
obtener representaciones a partir de cohomología geométrica.
Las hipótesis incluyen rango impar, condiciones sobre el cuerpo
y los modelos integrales, lugares no descompuestos y niveles
prescritos, una pareja auxiliar sobre $\ell$ y separación de pesos.

El resultado de esta etapa es una representación auxiliar con
ciertas ocurrencias de pesos y relaciones de anulación.
Para cada lugar adicional se obtiene una representación que
satisface la comparación requerida allí; la representación
auxiliar puede depender de ese lugar. Esta es la forma de
\texttt{HodgeTrAt}, ensamblada por \texttt{GeoOneOfParts}.
No proporciona todavía todas las multiplicidades ni el rango final.

El levantamiento theta permite reconocer las líneas adicionales.
Para un espacio hermitiano definido $X$ de dimensión $n$ y un espacio
antihermitiano definido $W$ de dimensión $d$ suficientemente grande,
la fuente fija retículos, matriz de Gram escalar y caracteres
de desdoblamiento con restricciones prescritas a $\mathbb A_F^\times$
y exponentes $m_X(v),m_W(v)$. Sea $U$ una representación automorfa
de $U(X)$ con peso dominante $\lambda$.
Aquí $d$ designa una dimensión, no el primo auxiliar de la
Nota~\ref{sec:auxiliares}. Si
\[
 d+m_W(v)\le2\lambda_k(v)
       \quad\text{para todos }v\mid\infty\text{ y }1\le k\le n,
\]
el peso del levantamiento viene dado por
\begin{equation}
 2\lambda'_k(v)=
 \begin{cases}
  2\lambda_k(v)-d-m_W(v)+n+m_X(v),&1\le k\le n,\\
  n+m_X(v),&n<k\le d.
 \end{cases}
 \label{eq:gal-theta-pesos}
\end{equation}
Los umbrales se toman después de fijar los datos pertinentes
y la conclusión se aplica a los $d$ admisibles suficientemente
grandes. No se deduce una fórmula universal $d_0=2n$.
Además de \eqref{eq:gal-theta-pesos}, se necesita el teorema de
no anulación con todos sus datos locales.

Un mecanismo elemental ayuda a leer esa no anulación.
En una representación unitaria de un grupo unimodular, para una
función de prueba $f$ que permita formar los operadores integrados,
$f^*(g)=\overline{f(g^{-1})}$ y
\[
 \langle\pi(f*f^*)v,v\rangle
       =\|\pi(f)^*v\|^2\ge0.
\]
La desigualdad es estricta cuando $\pi(f)^*v\ne0$.
Esta identidad no construye por sí sola el vector que debe sobrevivir.

\begin{lema}[Contribución compacta y cola]
\label{lem:gal-positividad}
Sea $kc$ integrable en un grupo medido y $K$ un subconjunto medible.
Si
\begin{equation}
 \int_K kc=A>0,\qquad
 \int_{G\setminus K}|kc|<A,
 \label{eq:gal-cola}
\end{equation}
entonces $\Ree\int_G kc>0$.
\end{lema}

\begin{proof}
La parte real de la integral en el complemento es al menos
$-\int_{G\setminus K}|kc|$. Al sumarle $A$ se obtiene un número
estrictamente positivo.
\end{proof}

Si $K$ es un subgrupo compacto, $c$ es el coeficiente matricial de un vector
unitario de una representación unitaria irreducible $\tau$ de $K$
y $k=\overline c$ en $K$, la ortogonalidad de Schur da
$A=\operatorname{vol}(K)/\dim\tau$.
Construir conjuntamente los núcleos, la identidad de producto interno,
la integrabilidad y la cota estricta de la cola es la entrada theta.
La desigualdad elemental no sustituye ese trabajo.

\subsection{De una resta virtual a un sumando real}

Sea $k$ un cuerpo algebraicamente cerrado de característica cero.
Cuando se hable de continuidad, $G$ será un grupo topológico,
$k$ un cuerpo topológico Hausdorff y $\{\pm1\}$ tendrá
la topología discreta. La aplicación aritmética usa
$k=\overline\Q_\ell$ con su topología $\ell$-ádica.
Para representaciones finitas semisimples, la clase $[R]$ registra
la multiplicidad de cada irreducible. En el grupo de Grothendieck
se permiten multiplicidades negativas; una resta de clases no
tiene por qué representar un espacio vectorial con acción.

\begin{proposicion}[Positividad de una diferencia común]
\label{prop:gal-diferencia}
Sean $R_1,R_2,A_1,A_2$ representaciones semisimples de un grupo.
Supongamos que $A_1$ y $A_2$ no comparten ningún constituyente
irreducible y que
\begin{equation}
 [R_1]-[A_1]=[R_2]-[A_2].
 \label{eq:gal-diferencia}
\end{equation}
Entonces existe una representación semisimple $V$, única hasta
isomorfía, tal que
\begin{equation}
 R_1\simeq V\oplus A_1,\qquad
 R_2\simeq V\oplus A_2.
 \label{eq:gal-sumandos}
\end{equation}
Si los datos son continuos, $V$ puede elegirse continua.
\end{proposicion}

\begin{proof}
Sea $m(S)$ la multiplicidad de un irreducible $S$ en la diferencia
común. Si $S$ no aparece en $A_1$, entonces $m(S)$ es su
multiplicidad en $R_1$ y es no negativa. Si aparece en $A_1$,
no aparece en $A_2$; usando la segunda expresión,
$m(S)$ es su multiplicidad en $R_2$, también no negativa.
Solo hay finitos $S$ con $m(S)\ne0$. La suma de $m(S)$ copias
de cada uno define $V$ y prueba las dos isomorfías.
Las multiplicidades determinan su clase de isomorfía.
Cada constituyente se realiza dentro de uno de los $R_i$;
las acciones restringidas y su suma finita son continuas.
\end{proof}

La separación necesaria es entre $A_1$ y $A_2$ como
representaciones, no entre sus nombres ni entre eigenvalores
evaluados en un único elemento. En la aplicación, son listas de
caracteres de Hecke cuya separación se conserva bajo realización
de Galois; esta última afirmación requiere el teorema de unicidad
de caracteres que usa \texttt{PCharInj}.

Las ecuaciones \eqref{eq:gal-sumandos} permiten retirar líneas
una a una mediante complementos invariantes.
El código \texttt{PPeelAll} desarrolla esa extracción.
Una desigualdad entre dimensiones totales no habría excluido
una multiplicidad negativa de un constituyente particular.

\subsection{De Frobenius a la igualdad de representaciones}
\label{subsec:gal-trazas}

Para una representación $R$ ponemos
$P_R(g,T)=\det(T-R(g))$. Las comparaciones theta dan, en los
elementos de Frobenius pertinentes, una identidad de la forma
\begin{equation}
 P_{R_1}(g,T)P_{A_2}(g,T)
     =P_{R_2}(g,T)P_{A_1}(g,T).
 \label{eq:gal-productos}
\end{equation}
Todos los factores son polinomios mónicos no nulos.
Chebotarev y la continuidad de cada coeficiente extienden esa
identidad a todo el grupo de Galois, una vez fijadas las
normalizaciones y las restricciones al mismo cuerpo.
Este uso del teorema de densidad es una entrada clásica;
\texttt{PDenseId} efectúa el enlace en la construcción.

\begin{lema}[Las trazas determinan una representación semisimple]
\label{lem:gal-trazas}
Dos representaciones finitas semisimples sobre $k$, con la misma
traza en todo elemento del grupo, son isomorfas.
\end{lema}

\begin{proof}
Sean $U,W$ las representaciones y $B$ el espacio generado por
la imagen diagonal del grupo en $\operatorname{End}(U\oplus W)$.
Es un álgebra finita dimensional: contiene la identidad y
el producto de dos matrices de la imagen vuelve a estar en ella.
Sus submódulos son exactamente los subespacios invariantes por
el grupo. Por tanto $U\oplus W$ es un $B$-módulo fiel semisimple.

El radical de $B$ aniquila cada módulo simple y, por fidelidad,
es cero. El teorema de estructura de álgebras semisimples escribe
$B$ como un producto de álgebras de matrices sobre $k$.
Si $S_j$ es el módulo simple de un factor y $e_j$ su idempotente
central, las trazas de $e_j$ en $U$ y $W$ son
$m_j\dim S_j$ y $n_j\dim S_j$. La igualdad de trazas en los
elementos del grupo se extiende por linealidad a $B$, así que
esos números coinciden. En característica cero podemos dividir
por $\dim S_j$ y obtenemos $m_j=n_j$ para todo $j$.
Esto da $U\simeq W$.
\end{proof}

El coeficiente de grado inmediatamente inferior al máximo en
\eqref{eq:gal-productos} da igualdad de trazas de
$R_1\oplus A_2$ y $R_2\oplus A_1$.
El lema produce la isomorfía; la
Proposición~\ref{prop:gal-diferencia} produce el sumando $V$.
Después se obtiene
\begin{equation}
 \dim V=\dim R_i-\dim A_i,\qquad
 P_{R_i}(g,T)=P_V(g,T)P_{A_i}(g,T).
 \label{eq:gal-cancelacion}
\end{equation}
La cancelación se hace en $k[T]$. No se divide por
$P_{A_i}(g,\alpha)$ en un eigenvalor $\alpha$,
que podría ser cero. Por ello se permiten eigenvalores
compartidos entre $V$ y $A_i$.

\subsection{Transporte local, pesos y giro final}

Una descomposición global $R=V\oplus A$ sigue siendo una
descomposición al restringir a un grupo de descomposición local.
No es necesario que esa restricción sea semisimple:
el proyector global sobre $V$ sigue siendo invariante localmente.
La no ramificación y las relaciones polinómicas de anulación
pasan directamente a $V$.

\begin{lema}[Defectos de dimensión y sumandos]
\label{lem:gal-periodos}
Sea $D$ un funtor aditivo de periodos para el que
$\dim D(U)\le\dim U$, sobre los cuerpos de comparación apropiados.
Si $R=V\oplus A$ y $\dim D(R)=\dim R$, entonces
$\dim D(V)=\dim V$ y $\dim D(A)=\dim A$.
\end{lema}

\begin{proof}
Por aditividad,
\[
 0=\dim R-\dim D(R)
   =(\dim V-\dim D(V))+(\dim A-\dim D(A)).
\]
Los dos términos son no negativos, luego ambos son cero.
\end{proof}

Las propiedades de comparación y la desigualdad de dimensión
de los funtores son entradas de teoría de Hodge $\ell$-ádica.
Una vez disponibles, el lema prueba la estabilidad por sumandos
de la propiedad de de Rham o cristalina.
Para la semiestabilidad potencial se toma primero una extensión
finita donde $R$ sea semiestable y se aplica allí el argumento.
Todas estas operaciones deben ocurrir en una misma realización
local compatible.

Las filtraciones y los operadores de Frobenius respetan la suma.
Si los pesos de $R$ son $J\sqcup J_A$ y los de $A$ son exactamente
$J_A$, entonces los de $V$ son $J$, con multiplicidades.
De modo semejante se cancelan factores mónicos en los polinomios
de inercia y de Frobenius. Si $N_R=0$, también $N_V=0$.
La compatibilidad del funtor de Weil--Deligne con esas comparaciones
es otra entrada; la cancelación posterior es álgebra lineal.
Estos pasos se separan en \texttt{PClSub}.

Queda un giro final. Para no confundirlo con el carácter cuadrático
CM $\kappa$ de las Notas~\ref{sec:periodos} y~\ref{sec:parametros},
lo escribimos $\kappa_{\mathrm{tw}}$.
Antes del giro, los eigenvalores deseados están multiplicados por
$\kappa_{\mathrm{tw}}^{-1}$ y los pesos desplazados en $m$.
El carácter de Galois asociado tiene peso $-m$. Por tanto
\begin{equation}
 (\alpha\,\kappa_{\mathrm{tw}}^{-1})
       \kappa_{\mathrm{tw}}=\alpha,\qquad
 (h+m)+(-m)=h.
 \label{eq:gal-giro}
\end{equation}
El mismo tensor corrige Frobenius y Hodge;
\texttt{PAtPlace} conserva ambas normalizaciones.

Finalmente, fijemos una extracción $V$. Para cada nuevo lugar $v$,
la entrada geométrica puede construir otra extracción $V_v$
con la propiedad deseada allí. Si ambas tienen los mismos
polinomios no ramificados fuera de un conjunto finito,
Chebotarev y el Lema~\ref{lem:gal-trazas} dan $V_v\simeq V$.
La propiedad local se transporta por esa isomorfía.
Así se obtiene una representación global fija, aunque los
objetos geométricos auxiliares varíen con $v$.

\subsection{Restricción y prolongación en índice dos}

\begin{lema}[Restricción a un subgrupo normal de índice finito]
\label{lem:gal-restriccion}
La restricción de una representación finita semisimple a un
subgrupo normal de índice finito es semisimple.
\end{lema}

\begin{proof}
Basta tratar una representación irreducible $W$ de $G$.
Si $H\triangleleft G$, un subespacio $H$-invariante no nulo de
dimensión mínima es un $H$-módulo simple $U$.
Por normalidad, cada trasladado $gU$ vuelve a ser simple para $H$.
La suma de esos trasladados, tomando representantes de $G/H$,
es no nula e invariante por $G$, luego es todo $W$.
Una suma finita de submódulos simples es semisimple:
al añadir uno, su intersección con la suma anterior es cero
o es todo él. Esto prueba la afirmación.
\end{proof}

\begin{lema}[Prolongaciones en índice dos]
\label{lem:gal-prolongacion}
Sea $H$ un subgrupo abierto normal de índice dos de $G$ y sea
$r$ una representación continua semisimple de $H$ sobre $k$.
Si $r$ es isomorfa a todos sus conjugados por $G$, admite una
prolongación continua semisimple a $G$.
Solo hay finitas clases de isomorfía de esas prolongaciones.
\end{lema}

\begin{proof}
Tomemos $t\in G\setminus H$. La conjugación por $t$ permuta los
constituyentes irreducibles de $r$ en órbitas de tamaño uno o dos;
sus multiplicidades son constantes en cada órbita.

Si $s$ es un constituyente fijo, existe un entrelazador invertible
$A$ con
\[
 A\,s(h)\,A^{-1}=s(tht^{-1})\qquad(h\in H).
\]
Por el lema de Schur, $s(t^2)^{-1}A^2=c\,1$ para un escalar
$c\ne0$. Como $k$ es algebraicamente cerrado, sustituimos $A$
por $c^{-1/2}A$. Ahora
\begin{equation}
 A^2=s(t^2).
 \label{eq:gal-prolongacion}
\end{equation}
Las relaciones de conjugación y \eqref{eq:gal-prolongacion}
permiten definir $s(t)=A$ y prolongar $s$ a $G$.
Las dos prolongaciones irreducibles se obtienen con $A$ y $-A$:
difieren por el carácter no trivial de $G/H$.

Si $s$ aparece con multiplicidad $m$, toda prolongación de ese
bloque isótipo tiene la forma anterior tensorizada con un operador
$J$ del espacio de multiplicidades, y $J^2=1$.
En característica cero $J$ es diagonalizable con eigenvalores
$1,-1$. Su clase queda determinada por
\begin{equation}
 m_++m_-=m,\qquad m_+,m_-\ge0.
 \label{eq:gal-elecciones}
\end{equation}
Hay exactamente $m+1$ posibilidades para ese bloque.

Si la órbita es $\{s,s^t\}$ con $s\not\simeq s^t$,
la representación $\operatorname{Ind}_H^G s$ es irreducible
y se restringe a $s\oplus s^t$.
La acción del coset no trivial intercambia los dos componentes;
cambios de base eliminan las elecciones restantes.
El bloque de multiplicidad $m$ se prolonga como $m$ copias
de esa representación inducida.

Tomamos la suma sobre las órbitas.
La descripción de los bloques agota las prolongaciones
semisimples y da finitas clases.
La acción es continua en $H$ y en su otro coset, donde es
el producto de una matriz fija y una función continua.
Como ambos cosets son abiertos, es continua en $G$.
\end{proof}

La prueba no supone que $G$ ni su imagen sean finitos.
La ambigüedad de signo es real: el carácter trivial y el carácter
no trivial de $G/H$ tienen la misma restricción a $H$.
Una sola restricción cuadrática no distingue todas las
prolongaciones.

\subsection{Por qué casi todos los núcleos conservan la imagen}

En este apartado la clausura es la \emph{clausura de Zariski}
de una imagen matricial: el conjunto de ceros de todas las
ecuaciones polinómicas que satisface esa imagen.
No es la clausura en la topología profinita de $G$.

\begin{lema}[Excepciones cuadráticas finitas]
\label{lem:gal-zariski}
Sea $R:G\to\operatorname{GL}(U)$ una representación finita
sobre $k$. Salvo finitos caracteres cuadráticos
$\eta:G\to\{\pm1\}$, se tiene
\begin{equation}
 \overline{R(\ker\eta)}^{\,\mathrm{Zar}}
       =\overline{R(G)}^{\,\mathrm{Zar}}.
 \label{eq:gal-zariski}
\end{equation}
Por tanto, para una familia infinita de caracteres distintos,
la restricción a casi todos sus núcleos conserva los subespacios
invariantes y los entrelazadores de cualquier lista finita
de representaciones fijada de antemano.
\end{lema}

\begin{proof}
Sea $\mathcal G=\overline{R(G)}^{\,\mathrm{Zar}}$,
$\mathcal G^\circ$ su componente de la identidad y
$M=R^{-1}(\mathcal G^\circ)$.
El grupo algebraico $\mathcal G$ tiene finitas componentes;
por eso $M$ es normal de índice finito en $G$.
No se necesita afirmar que $M$ sea abierto en la topología original.

Si $\eta$ no es trivial en $M$ y $H=\ker\eta$,
entonces $G=HM$.
La clausura $\mathcal H=\overline{R(H)}^{\,\mathrm{Zar}}$
tiene índice a lo sumo dos en $\mathcal G$ y contiene
$\mathcal G^\circ$. Además, $G=HM$ implica que
$\mathcal H$ alcanza todas las componentes de $\mathcal G$.
Así $\mathcal H=\mathcal G$.
Los caracteres que podrían fallar son triviales en $M$,
luego factorizan por el grupo finito $G/M$;
solo hay finitos.

Para una lista $R_1,\ldots,R_s$, aplicamos el argumento
a $R_1\oplus\cdots\oplus R_s$.
La invariancia de un subespacio y las ecuaciones
$T R_j(g)=R_k(g)T$ son condiciones polinómicas en las matrices.
Si valen en la imagen del núcleo, valen en su clausura y,
por \eqref{eq:gal-zariski}, en toda la imagen de $G$.
\end{proof}

El mismo argumento se aplica dentro de un subgrupo $H_0$
de índice dos. Además, la aplicación de restricción
\begin{equation}
 \eta\longmapsto\eta|_{H_0}
 \label{eq:gal-fibras}
\end{equation}
tiene fibras de tamaño a lo sumo dos: el cociente de dos
caracteres con la misma restricción factoriza por $G/H_0$.
Por tanto una familia infinita de caracteres distintos
sigue dando infinitos caracteres distintos tras restringir,
y excluir finitos caracteres restringidos excluye solo
finitos de la familia inicial.

\subsection{El pegado cuadrático, con sus elecciones finitas}
\label{subsec:gal-pegado}

\begin{teorema}[Pegado cuadrático de A93]
\label{teo:gal-pegado}
Sea $\{\eta_i\}$ una familia infinita de caracteres cuadráticos
continuos distintos de $G$, y $H_i=\ker\eta_i$.
Para cada $i$, sea $r_i:H_i\to\operatorname{GL}_a(k)$
continua semisimple. Supongamos:
\begin{enumerate}
 \item la traza de $r_i$ es invariante por conjugación por $G$;
 \item para todo par $i,j$, las trazas de $r_i$ y $r_j$
       coinciden en $H_i\cap H_j$.
\end{enumerate}
Existe una representación continua semisimple
$\rho:G\to\operatorname{GL}_a(k)$ tal que
\begin{equation}
 \rho|_{H_i}\simeq r_i
       \quad\text{para todos salvo finitos }i.
 \label{eq:gal-eventual}
\end{equation}
\end{teorema}

\begin{proof}
Retiramos el posible carácter trivial y fijamos un índice $0$.
La invariancia de trazas, la semisimplicidad y el
Lema~\ref{lem:gal-trazas} dan invariancia por conjugación
hasta isomorfía.
El Lema~\ref{lem:gal-prolongacion} proporciona una lista finita
$R_1,\ldots,R_s$ de todas las prolongaciones semisimples de $r_0$.

Pongamos $L_i=H_0\cap H_i$.
Aplicando el Lema~\ref{lem:gal-zariski} a $r_0$ en $H_0$
y usando \eqref{eq:gal-fibras}, vemos que, salvo finitos $i$,
la restricción a $L_i$ conserva la irreducibilidad y la distinción
de los constituyentes de $r_0$.
Por el Lema~\ref{lem:gal-restriccion}, las restricciones
$r_i|_{L_i}$ y $r_0|_{L_i}$ son semisimples.
Sus trazas coinciden, luego
\begin{equation}
 r_i|_{L_i}\simeq r_0|_{L_i}.
 \label{eq:gal-interseccion}
\end{equation}

Expliquemos por qué esto hace de $r_i$ la restricción de algún
$R_j$, que es el punto finito de la prueba.
Para $i\ne0$, el mapa $H_i/L_i\to G/H_0$ es un isomorfismo.
Podemos elegir el representante del coset no trivial de $G/H_0$
dentro de $H_i$.
Su acción sobre los constituyentes restringidos de $r_0$
tiene las mismas órbitas de tamaño uno o dos que antes:
la restricción ha conservado irreducibilidad y distinción.
En cada órbita fija, las dos prolongaciones a $H_i$
difieren por $\eta_0|_{H_i}$; son las restricciones de las
dos prolongaciones a $G$.
Si la multiplicidad es $m$, las opciones
$m_++m_-=m$ también coinciden.
En una órbita de tamaño dos, la única prolongación inducida
con esa multiplicidad se obtiene igualmente por restricción.
En efecto, $G=H_iH_0$ da
\[
 (\operatorname{Ind}_{H_0}^G s)|_{H_i}
   \simeq\operatorname{Ind}_{L_i}^{H_i}(s|_{L_i}).
\]
El mapa entre los modelos tensoriales envía $h\otimes v$ a
$h\otimes v$; es sobreyectivo por $G=H_iH_0$ y ambos espacios
tienen dimensión $2\dim s$, luego es un isomorfismo.
Así todas las elecciones que describen $r_i$ en
\eqref{eq:gal-interseccion} aparecen en la lista $R_1,\ldots,R_s$.

Apliquemos ahora el Lema~\ref{lem:gal-zariski} a
$R=R_1\oplus\cdots\oplus R_s$.
Elegimos $i_1$ fuera de las excepciones anteriores y de las
excepciones de este lema. La imagen $R(H_{i_1})$ tiene entonces
la misma clausura que $R(G)$.
Tomamos $\rho$ en la lista de modo que
$\rho|_{H_{i_1}}\simeq r_{i_1}$.

Aplicando de nuevo el lema dentro de $H_{i_1}$,
para todos salvo finitos $i$ la imagen
$R(H_i\cap H_{i_1})$ tiene esa misma clausura.
Ya sabemos que $r_i\simeq R_j|_{H_i}$ para algún $j$.
La compatibilidad de trazas y la semisimplicidad dan una
isomorfía entre $R_j$ y $\rho$ en $H_i\cap H_{i_1}$.
Las ecuaciones de entrelazamiento se extienden por densidad
de Zariski a $R(G)$, luego $R_j\simeq\rho$ en $G$.
Esto prueba \eqref{eq:gal-eventual}.
La continuidad y la semisimplicidad de $\rho$ provienen
de la lista de prolongaciones.
\end{proof}

Esta reconstrucción sigue la prueba expositiva de A93,
Proposición~10.5, pp.~82--83; A presenta el núcleo en
Proposición~9.4, p.~45.
La declaración \texttt{IndexTwoPatch.main} del código llega
a igualdad eventual de trazas mediante finitud de caracteres
cuadráticos en una representación adjunta y control de
entrelazadores torcidos.
\texttt{ResInjective} convierte esas trazas en isomorfías
mediante semisimplicidad.
Las dos rutas resuelven el mismo descenso, con pruebas distintas.

\subsection{Una ampliación: recuperación de toda la familia}
\label{subsec:gal-todos}

La conclusión eventual basta para la aplicación de A93.
Bajo sus hipótesis abstractas literales, se puede reforzar.
Conservamos aquí la cuantificación sobre \emph{todo par}
de índices; debilitarla cambiaría el argumento.

\begin{proposicion}[Ampliación del pegado]
\label{prop:gal-todos}
Bajo las hipótesis del Teorema~\ref{teo:gal-pegado},
la representación $\rho$ obtenida satisface
\begin{equation}
 \rho|_{H_i}\simeq r_i
                \quad\text{para todo }i.
 \label{eq:gal-todos}
\end{equation}
Además es única hasta isomorfía entre las representaciones
semisimples con las restricciones eventuales prescritas.
\end{proposicion}

\begin{proof}
Fijemos un índice $i$. Consideremos en $H_i$ la representación
$U=(\rho|_{H_i})\oplus r_i$.
Si $H_i$ tiene índice dos, las restricciones $\eta_j|_{H_i}$
tienen fibras de tamaño a lo sumo dos; si $H_i=G$, son ya
distintas. En ambos casos siguen siendo infinitas.
El Lema~\ref{lem:gal-zariski} permite elegir $j$ fuera del
conjunto excepcional de \eqref{eq:gal-eventual}, de modo que
\[
 \overline{U(H_i\cap H_j)}^{\,\mathrm{Zar}}
       =\overline{U(H_i)}^{\,\mathrm{Zar}}.
\]
En la intersección, la traza de $\rho$ coincide con la de $r_j$,
y esta coincide con la de $r_i$ por compatibilidad.
La diferencia de esas trazas es una función polinómica de las
matrices de $U$, luego se anula en todo $H_i$.
El Lema~\ref{lem:gal-restriccion} da semisimplicidad de
$\rho|_{H_i}$; el Lema~\ref{lem:gal-trazas} da
\eqref{eq:gal-todos} para ese $i$.

Para la unicidad, sean $\rho,\rho'$ dos soluciones eventuales.
Aplicamos el Lema~\ref{lem:gal-zariski} a $\rho\oplus\rho'$.
Elegimos $i$ fuera de sus excepciones y de las de ambas
soluciones. La isomorfía entre sus restricciones a $H_i$
es un entrelazador invertible cuyas ecuaciones valen en una
imagen densa de Zariski; valen, por tanto, en toda la imagen
de $G$. Así $\rho\simeq\rho'$.
\end{proof}

Esta mejora es una consecuencia matemática añadida al cuaderno.
No se atribuye como enunciado a \texttt{IndexTwoPatch.main}
ni se presenta como un nuevo archivo verificado por Lean.
Tampoco corrige un error del enunciado más débil de la fuente:
explica una consecuencia adicional de sus hipótesis completas.

\subsection{Recuperar propiedades en un lugar que se descompone}

Para aplicar el pegado a $G_E$ se usan extensiones cuadráticas
$E_i/E$ y $H_i=G_{E_i}$. La invariancia y la compatibilidad
en $H_i\cap H_j=G_{E_iE_j}$ provienen de cambio de base
automorfo y de los polinomios de Frobenius.
Chebotarev los extiende a todo el grupo de la intersección.
La mera existencia de muchas representaciones del mismo rango
no proporciona esas igualdades.

\begin{proposicion}[Transporte por descomposición local]
\label{prop:gal-split}
Supongamos que $\rho|_{G_{E_i}}\simeq r_i$ y que un lugar $v$
de $E$ se descompone completamente en $E_i$.
Para un lugar $u$ sobre $v$, toda propiedad de $r_i$ en $u$
invariante por isomorfía de representaciones locales se
transfiere a $\rho$ en $v$, con los datos de comparación
compatibles.
\end{proposicion}

\begin{proof}
La descomposición da $E_{i,u}\simeq E_v$.
Con elecciones compatibles de clausuras algebraicas y lugares,
el grupo de descomposición original $D_v$ está contenido
en $G_{E_i}$ y es el mismo grupo local que se utiliza en $u$.
La isomorfía global restringida identifica las dos
representaciones de ese grupo.
Transportamos por ella los operadores, las filtraciones y los
puentes de comparación que intervengan en la propiedad.
No se ha cambiado el cuerpo local.
\end{proof}

En la ruta eventual de A93, para cada $v$ se construyen
infinitas extensiones permitidas donde $v$ se descompone.
Ello permite elegir una fuera del conjunto finito excepcional.
Los cuantificadores tienen el orden
\begin{equation}
 \exists\rho\ \exists I_{\mathrm{exc}}\text{ finito}\
 \forall v\ \exists i_v\notin I_{\mathrm{exc}}
 \quad
 \left[
  \begin{array}{l}
   v\text{ se descompone en }E_{i_v},\\
   \rho|_{G_{E_{i_v}}}\simeq r_{i_v}.
  \end{array}
 \right]
 \label{eq:gal-split-orden}
\end{equation}
El índice $i_v$ puede depender de $v$; $\rho$ ya está fijada.
La infinitud por lugar es suficiente para esta ruta de la fuente.
La Proposición~\ref{prop:gal-todos} refuerza el pegado abstracto,
pero no construye por sí misma extensiones aritméticas con
descomposición y con todas las demás condiciones exigidas.

La descomposición local tiene una función indispensable en este
transporte directo. Para una extensión cuadrática local ramificada
$L/K$, el carácter cuadrático
$\eta_{L/K}:G_K\to\{\pm1\}$ es ramificado en $K$, pero su
restricción a $G_L$ es trivial.
Por tanto la no ramificación tras una extensión local arbitraria
no demuestra no ramificación sobre $K$.
En un lugar descompuesto no se introduce esa extensión local.

El pegado por extensiones y el uso de familias suficientemente
generales tienen precedentes explícitos en
\cite{sorensen-patching}. Se utiliza aquí como antecedente;
las hipótesis de A93 se enuncian y comprueban en su propio marco.
Los dos usos de densidad también se mantienen separados:
Chebotarev extiende igualdades de Frobenius a grupos de Galois,
y la clausura de Zariski conserva entrelazadores de una lista
finita de representaciones.
Un subgrupo abierto de índice dos no es denso en el grupo
profinito original.

\subsection{Regreso a los periodos y a la cohomología}

La entrada de Galois debe aplicarse a la representación automorfa
que surge de los mismos datos de las
Notas~\ref{sec:periodos}--\ref{sec:parametros}.
El giro \eqref{eq:gal-giro}, la igualdad de pesos con multiplicidad
y el alcance exacto de las dos ramas locales evitan sustituir
esa representación por otra con datos normalizados de forma distinta.

En la Nota~\ref{sec:cohomologia} se usa Frobenius geométrico.
Para el carácter ciclotómico $\chi_\ell$, fuera de $\ell$,
\begin{equation}
 \chi_\ell(\operatorname{Frob}^{\mathrm{geom}}_v)=q_v^{-1},
 \qquad
 \chi_\ell^{-e}(\operatorname{Frob}^{\mathrm{geom}}_v)=q_v^e.
 \label{eq:gal-ciclotomico}
\end{equation}
Esto concuerda con peso de Hodge $-1$ para $\chi_\ell$
y con el convenio de torcidos de la nota anterior.
El cotejo de este bloque no exige cambiar esas ecuaciones.
Sí hace explícito por qué la monodromía cero, los pesos y las
hipótesis sobre el cuerpo local deben acompañar a la comparación.

Hay un límite adicional: una representación semisimple sobre
$\overline\Q_\ell$ no especifica un retículo integral ni una
clase de extensión. Por ejemplo, el grupo aditivo $\Z_\ell$
actúa continuamente mediante
\begin{equation}
 \rho(z)=
 \begin{pmatrix}1&z\\0&1\end{pmatrix}.
 \label{eq:gal-unipotente}
\end{equation}
Su semisimplificación es la representación trivial de dimensión
dos. Ambas tienen traza $2$ y polinomio $(T-1)^2$ para todo $z$,
pero no son isomorfas: $\rho(1)\ne1$.
La información que detecta la extensión se pierde al semisimplificar.

Por eso la entrada \texttt{Row0b} no reemplaza los teoremas
integrales que enlazan congruencias de Hecke, extensiones y
condiciones locales de Selmer.
Las longitudes y los denominadores de la
Nota~\ref{sec:cohomologia} requieren precisamente esa información.
Del mismo modo, el pegado no produce las constantes uniformes
de las ventanas de la Nota~\ref{sec:parametros}
ni las estimaciones analíticas de la Nota~\ref{sec:auxiliares}.
Aporta la representación y las comparaciones que esos pasos consumen.

\subsection{Concordancia localizada con Lean}
\label{subsec:gal-lean}

El cotejo conserva la revisión \texttt{7cbf7a0} de [L].
Separa las declaraciones que expresan una propiedad de los
teoremas que la producen, y conserva las dependencias profundas
como tales. El manifiesto noveno registra rutas completas,
hashes y rangos exactos de lectura.
Comprende 36 archivos: 32 lecturas completas y cuatro selectivas,
con 3299 líneas leídas. Los archivos pueden coincidir con otros
cotejos anteriores; el recuento describe este tramo.

\begin{center}
\small
\begin{tabularx}{\textwidth}{@{}p{0.27\textwidth}YY@{}}
\toprule
Paso & Archivo o declaración & Alcance observado\\
\midrule
Polarización y salida &
\texttt{DualIsConjByOne}; \texttt{WeightsConjBis}; \texttt{RowSrc} &
Hipótesis de signo, pesos y ramas locales; productor que ensambla la salida.\\
Realizaciones locales &
\texttt{exists\_locPrime\_tied}; \texttt{exists\_locBridge} &
Testigos para datos de periodos y para la valoración que induce el puente.\\
Geometría y theta &
\texttt{HodgeTrAt}; \texttt{GeoOneOfParts};
las dos rutas de \texttt{TLiftK} &
Objetos auxiliares, dependencias geométricas y elección por lugar.\\
Extracción &
\texttt{PDenseId}; \texttt{PCharInj}; \texttt{PPeelAll} &
Identidad densa, separación real de caracteres y retirada sucesiva de líneas.\\
Transporte local &
\texttt{PClSub}; \texttt{PAtPlace} &
Sumandos, periodos, cancelación polinómica y giro final.\\
Pegado &
\texttt{EqOneOrEqOfKer}; \texttt{ResInjective} &
Trazas eventuales por caracteres en la adjunta; isomorfías tras semisimplicidad.\\
Descenso por lugar &
\texttt{CoreAllOfRowStarCoreInp} &
Selección fuera de finitos y control del grupo de descomposición.\\
Asignación auxiliar &
\texttt{IsRecTOfIsRec}; \texttt{RecTExists} &
Leyes globales débiles y leyes fuertes calificadas por la clase buena.\\
\bottomrule
\end{tabularx}
\end{center}

Los nombres de archivo abreviados no sustituyen sus rutas.
\href{https://github.com/lachenmann/riemann-notas/blob/main/fuentes/README.md}{El inventario de fuentes}
las proporciona, junto con enlaces fijados a la revisión de código.
En \texttt{RowSrc}, el teorema
\texttt{galoisLocalGlobalStrong\_recParameter}, del espacio de nombres
\texttt{UpperTheory.Gal.Int}, es un productor que invoca un ensamblaje theta;
leer su prueba breve no equivale a reconstruir todos sus teoremas
importados. También se leyeron productores de las interfaces
geométrica, theta, de extracción y de asignación local.

La diferencia entre la prueba de pegado por Zariski y la prueba
del código por caracteres de la adjunta se deja registrada.
La Proposición~\ref{prop:gal-todos} es un desarrollo matemático
de estas notas. No se ejecutó Lean, no se obtuvo un informe de
axiomas y no se auditó la clausura de todas las dependencias.
Los hashes identifican lo leído; por sí solos no validan
la demostración global.

\subsection{Ejercicios resueltos}

\begin{ejercicio}\label{ej:gal-virtual}
Sean $S,T$ irreducibles no isomorfos y tómese
$R_1=R_2=T$, $A_1=A_2=S$.
Compruebe la igualdad virtual y explique por qué la conclusión
de la Proposición~\ref{prop:gal-diferencia} puede fallar.
\end{ejercicio}

\paragraph{Solución del Ejercicio~\ref{ej:gal-virtual}.}
Las dos diferencias son $[T]-[S]$.
La multiplicidad de $S$ es $-1$, así que no representa una
representación semisimple genuina.
Falla precisamente la separación de los constituyentes de
$A_1$ y $A_2$.
Incluso si $\dim T\ge\dim S$, la dimensión virtual no negativa
no elimina esa multiplicidad negativa.

\begin{ejercicio}\label{ej:gal-anulador}
Para $u\ne v$, calcule las trazas de los operadores diagonales
con multiconjuntos $\{u,u,v\}$ y $\{u,v,v\}$.
Explique por qué su aniquilador común y su dimensión no
determinan las multiplicidades.
\end{ejercicio}

\paragraph{Solución del Ejercicio~\ref{ej:gal-anulador}.}
Las trazas son $2u+v$ y $u+2v$, con diferencia $u-v\ne0$.
Ambos operadores tienen dimensión tres y son aniquilados
por $(T-u)(T-v)$, pero sus polinomios característicos son
los de \eqref{eq:gal-multiplicidades}.
El aniquilador solo obliga a que cada eigenvalor pertenezca
a $\{u,v\}$; no cuenta cuántas veces aparece.

\begin{ejercicio}\label{ej:gal-prolongaciones}
Suponga que un constituyente irreducible $s$ de $H$ es fijo
por conjugación y aparece con multiplicidad tres.
Cuente las prolongaciones de ese bloque en índice dos.
Haga lo mismo para una órbita $\{t,t^g\}$ no fija que aparece
con multiplicidad tres.
\end{ejercicio}

\paragraph{Solución del Ejercicio~\ref{ej:gal-prolongaciones}.}
En el primer caso los pares $(m_+,m_-)$ son
$(3,0),(2,1),(1,2),(0,3)$: hay cuatro clases.
No hay ocho clases, pues permutar tres copias indistinguibles
no produce una nueva representación.
En el segundo caso solo hay una: tres copias de
$\operatorname{Ind}_H^G t$.
Si ambos bloques aparecen a la vez y no hay otros,
hay cuatro prolongaciones en total.

\begin{ejercicio}\label{ej:gal-densidad}
En $G=\Z_2^2$, tome
$\eta(x,y)=(-1)^{x\bmod2}$ y
$R(x,y)=\begin{pmatrix}1&x\\0&1\end{pmatrix}$ sobre
$\overline\Q_2$.
Compare la clausura profinita de $\ker\eta$ con las clausuras
de Zariski de $R(G)$ y $R(\ker\eta)$.
\end{ejercicio}

\paragraph{Solución del Ejercicio~\ref{ej:gal-densidad}.}
El núcleo es $2\Z_2\times\Z_2$, cerrado y propio de índice dos;
su clausura profinita no es $G$.
Las imágenes tienen parámetros $x\in\Z_2$ y $x\in2\Z_2$,
respectivamente. Ambos conjuntos son infinitos.
Un polinomio de una variable sobre un cuerpo que se anula
en cualquiera de ellos debe ser cero.
Por ello ambas clausuras de Zariski son el mismo grupo
unipotente $\{\begin{pmatrix}1&t\\0&1\end{pmatrix}:t\in\overline\Q_2\}$.
Este ejemplo ilustra las dos topologías; no pretende por sí
solo proporcionar una familia infinita de caracteres cuadráticos.

\begin{ejercicio}\label{ej:gal-local}
Sea $L/K$ una extensión cuadrática local ramificada.
Muestre que $\eta_{L/K}|_{G_L}$ es no ramificada
aunque $\eta_{L/K}$ no lo sea.
Explique qué impide este ejemplo en el descenso local.
\end{ejercicio}

\paragraph{Solución del Ejercicio~\ref{ej:gal-local}.}
El núcleo de $\eta_{L/K}$ es $G_L$, de modo que la restricción
es trivial. La imagen de la inercia de $K$ en
$\operatorname{Gal}(L/K)$ es no trivial porque la extensión
es ramificada; por eso el carácter original es ramificado.
La trivialidad después de restringir no detecta la inercia
perdida al pasar a $L$.
En la Proposición~\ref{prop:gal-split}, en cambio,
$E_{i,u}=E_v$ y el grupo local no disminuye.

\subsection{Lo demostrado y el siguiente enlace}

Se han demostrado la extracción positiva, la determinación por
trazas, el argumento de periodos por sumandos, la finitud de
prolongaciones en índice dos, el pegado por clausuras de Zariski,
su ampliación a toda la familia y el transporte en lugares
descompuestos. También se han precisado la multiplicidad mansa,
la monodromía y el orden de las elecciones locales.

Quedan recibidas como entradas la construcción geométrica
con sus condiciones integrales, la no anulación theta uniforme,
el cambio de base automorfo, las leyes de los funtores de periodos
y la asignación local en el rango requerido.
La lectura localizada de sus productores no desarrolla todas
esas pruebas.

La Nota~\ref{sec:congruencias-selmer} desarrolla el siguiente enlace
mediante un modelo de congruencias, retículos y extensiones en
dimensión dos. Demuestra que una clase residual no trivial puede
anularse al pasar al módulo divisible con la acción levantada.
La comparación con Selmer conserva como hipótesis la existencia
de un mapa y el control de su núcleo y de las condiciones locales.
El paso aritmético completo remite a la Nota~\ref{sec:cohomologia},
A55, §8, pp.~35--42, y A93, §9, pp.~61--77;
no se deduce de la sola existencia de $\rho_{\pi,\ell}$.


